\chapter{The Unit Circle and General Angles}
\label{ch:ch04-unit-circle}

\section{Beyond Right Triangles}

The trigonometric functions were introduced in the previous chapter as
ratios associated with right triangles, and that approach provided a
powerful method for relating angles and side lengths. It also imposed a
limitation, since a right triangle can contain only acute angles: the angle
$\theta$ under discussion had to satisfy $0^\circ < \theta < 90^\circ$. What,
then, is the sine of $150^\circ$, or the cosine of $-45^\circ$, or the
meaning of an angle of $720^\circ$? None of these questions can be answered
using right triangles alone, and the difficulty is not with trigonometry
itself but with having tied the trigonometric functions to the interior
angle of a triangle in the first place. To remove the restriction, angle
must be redefined not as part of a triangle but as a rotation, and the
geometric object that makes this possible is the unit circle.

\section{The Unit Circle}

Recall the distance formula of Chapter~2: a point $(x,y)$ lies at distance
$\sqrt{x^2+y^2}$ from the origin. Requiring every point under consideration
to lie exactly one unit from the origin means requiring $\sqrt{x^2+y^2}=1$,
or equivalently $x^2+y^2=1$.

\begin{definition}
The \emph{unit circle} is the set of all points $(x,y)$ satisfying
$x^2+y^2=1$.
\end{definition}

The unit circle is not the entire plane; it is a constrained subset of it,
in exactly the sense of Chapter~0. Among all possible points $(x,y)$, only
those satisfying $x^2+y^2=1$ are admissible, and this single constraint
fixes magnitude at $1$ while leaving direction free to vary. The circle is
therefore precisely the geometric setting needed to study direction in
isolation, without the interference of size that occupied Chapter~1.

\section{Angles as Rotations}

\begin{definition}
An \emph{angle} is a rotation from the positive $x$-axis, called the
\emph{initial side}, to a rotated ray called the \emph{terminal side}. A
counterclockwise rotation is considered positive, and a clockwise rotation
is considered negative. The point where the terminal side meets the unit
circle is the point corresponding to that angle.
\end{definition}

This definition allows angles such as $120^\circ$, $-45^\circ$, and
$720^\circ$ to be spoken of naturally, since none of them is required to fit
inside a triangle: an angle is now simply an amount of rotation, unrestricted
in sign or size.

\section{Degree Measure and Radian Measure}

Angles have historically been measured in degrees, with a complete
revolution comprising $360^\circ$. Mathematics also makes frequent use of a
second unit, the radian.

\begin{definition}
One \emph{radian} is the angle that subtends an arc of the unit circle whose
length equals the circle's radius.
\end{definition}

Because the circumference of a circle of radius $r$ is $2\pi r$, a complete
revolution subtends an arc of length $2\pi r$, and dividing by the radius
shows that a complete revolution corresponds to an angle of $2\pi$ radians.
Hence $360^\circ = 2\pi$ radians, from which the conversion formulas
\[
\theta_{\text{rad}} = \theta_{\text{deg}}\cdot\frac{\pi}{180}, \qquad
\theta_{\text{deg}} = \theta_{\text{rad}}\cdot\frac{180}{\pi}
\]
follow immediately. Degrees and radians are two units for the same
geometric quantity, and this book will move between them as convenience
dictates.

\section{Extending Sine and Cosine}

Let $\theta$ be any angle, and rotate from the positive $x$-axis until
reaching the terminal side of $\theta$. Suppose this terminal side meets the
unit circle at the point $P = (x,y)$.

\begin{definition}
For any angle $\theta$, $\cos\theta = x$ and $\sin\theta = y$, where $(x,y)$
is the point at which the terminal side of $\theta$ meets the unit circle.
In words, cosine is the horizontal coordinate of this point and sine is its
vertical coordinate.
\end{definition}

This definition places no restriction on $\theta$ whatsoever, and applies
equally well to angles of any size or sign. It remains to check that nothing
established in Chapter~3 has been lost in the generalization.

\begin{theorem}
For an acute angle $\theta$ of a right triangle, the unit-circle definitions
of sine and cosine agree with the definitions of Chapter~3.
\end{theorem}

\begin{proof}
Let the right triangle have hypotenuse $r$, so that by Chapter~3,
$\cos\theta = \text{adjacent}/r$ and $\sin\theta = \text{opposite}/r$.
Dividing every side length of the triangle by $r$ scales the triangle onto
the unit circle, by the scalar multiplication of Chapter~2 with $k = 1/r$,
and the vertex corresponding to the angle $\theta$ moves to the point
$\left(\text{adjacent}/r,\ \text{opposite}/r\right)$. This is exactly the
point at which the terminal side of $\theta$ meets the unit circle, so
$x = \cos\theta$ and $y = \sin\theta$ under the new definition agree with
the ratios of the old one.
\end{proof}

Nothing from Chapter~3 has therefore been replaced; the unit-circle
definitions merely extend the right-triangle definitions to angles the
triangle could never have contained.

\section{Tangent and the Remaining Functions}

Since $\sin\theta = y$ and $\cos\theta = x$ on the unit circle, tangent is
defined as before by
\[
\tan\theta = \frac{\sin\theta}{\cos\theta} = \frac{y}{x}.
\]
Whenever $x=0$ this ratio is undefined, so tangent fails to exist at
$90^\circ$, $270^\circ$, and every angle coterminal with them. The
reciprocal functions are defined exactly as in Chapter~3:
\[
\csc\theta = \frac{1}{\sin\theta}, \qquad
\sec\theta = \frac{1}{\cos\theta}, \qquad
\cot\theta = \frac{1}{\tan\theta}.
\]

\section{Coterminal Angles}

Beginning at the positive $x$-axis and rotating through one complete
revolution returns to the same point of the unit circle, so adding
$360^\circ$ to an angle changes its description but not the point it names.

\begin{theorem}
For every integer $k$, the angles $\theta$ and $\theta + 360^\circ k$
determine the same point of the unit circle; equivalently, $\theta$ and
$\theta + 2\pi k$ determine the same point when measured in radians. Angles
related in this way are called \emph{coterminal}.
\end{theorem}

This is the first appearance, in the context of angle itself, of a theme
that will recur throughout the rest of the book: periodicity. The circle
repeats, and its geometry is inherently cyclic, so that any function defined
in terms of position on the circle must inherit that same repetition.

\section{Reference Angles}

Although the trigonometric functions are now defined for every angle,
evaluating them typically reduces to a much simpler problem, since every
angle can be associated with an acute angle in the first quadrant called its
\emph{reference angle}.

\begin{algorithmproc}{How to Evaluate a Trigonometric Function Using a Reference Angle}
Given an angle $\theta$, first locate the quadrant containing its terminal
side. Next determine the acute reference angle formed between the terminal
side and the $x$-axis. Evaluate the corresponding trigonometric function at
this first-quadrant reference angle, and finally apply the sign appropriate
to the original quadrant, using the fact that $\cos\theta$ is the horizontal
coordinate and $\sin\theta$ is the vertical coordinate of a point on the
unit circle.
\end{algorithmproc}

\begin{table}[h]
\centering
\begin{tabular}{lccc}
\toprule
Quadrant & $\sin\theta$ & $\cos\theta$ & $\tan\theta$ \\
\midrule
I   & $+$ & $+$ & $+$ \\
II  & $+$ & $-$ & $-$ \\
III & $-$ & $-$ & $+$ \\
IV  & $-$ & $+$ & $-$ \\
\bottomrule
\end{tabular}
\caption{Sign of each trigonometric function by quadrant.}
\end{table}

\begin{example}
Evaluate $\sin(210^\circ)$. Since $210^\circ = 180^\circ + 30^\circ$, the
angle lies in Quadrant III with reference angle $30^\circ$. From
Chapter~3, $\sin(30^\circ) = \tfrac12$, and sine is negative in Quadrant
III, so
\[
\sin(210^\circ) = -\frac12.
\]
\end{example}

\begin{practicenow}
Evaluate $\cos(135^\circ)$ using a reference angle.
\end{practicenow}

\begin{sidebar}{The Sexagenary Cycle and Simultaneous Periods}
A sharper illustration of periodicity than the calendar in general comes
from the traditional Chinese sexagenary cycle, which names successive
years, and historically days, by pairing one of ten symbols called
heavenly stems with one of twelve symbols called earthly branches. Both
sequences advance together, one position at a time, so that a given
pairing recurs only once both cycles have returned to their starting
point simultaneously. Since the stem cycle has length ten and the branch
cycle has length twelve, the combined cycle has length
$\operatorname{lcm}(10,12) = 60$, not the $10\times12 = 120$ pairings that
might naively be expected, because the two cycles advance in lockstep
rather than independently: only those pairings whose position differences
are consistent with a single shared step actually occur, cutting the
apparent count in half. This is ordinary modular arithmetic operating
inside a human institution centuries before the vocabulary of cyclic
groups existed to describe it. The unit circle exhibits the same
structure in continuous form: two angles are coterminal exactly when
their difference is a multiple of the circle's own period, $2\pi$, and
the interaction of two distinct periods, as will reappear when several
oscillations of different frequencies are added together later in this
book, likewise produces a combined period governed by a common multiple
rather than a simple product of the two.
\end{sidebar}

\section{Exact Values on the Unit Circle}

The special angles derived in Chapter~3 now acquire explicit geometric
locations on the circle:
\[
30^\circ \mapsto \left(\frac{\sqrt3}{2}, \frac12\right), \qquad
45^\circ \mapsto \left(\frac{\sqrt2}{2}, \frac{\sqrt2}{2}\right), \qquad
60^\circ \mapsto \left(\frac12, \frac{\sqrt3}{2}\right).
\]
Combined with the sign pattern of the reference-angle table, these three
points generate the exact trigonometric value of every multiple of
$30^\circ$ and $45^\circ$ around the entire circle, so that the unit circle
functions as a complete map of the exact values a course in trigonometry is
expected to know without a calculator.

\section{Applications}

The unit circle appears wherever a rotational phenomenon is studied: in the
motion of wheels and gears, in the orbits of planets, in rotating
machinery, and in the description of waves and oscillations. By fixing
magnitude at $1$ and allowing only direction to vary, the unit circle
supplies a single, universal geometric language in which every one of these
rotational systems can be described.

\section{Looking Ahead}

The unit circle allows the trigonometric functions to be defined for every
real angle, and it has revealed periodicity, symmetry, and the exact
geometric values underlying the special angles of Chapter~3. Everything
done in this chapter, however, remains static: an angle is specified, a
point is identified, a value is computed, and the process stops there. What
happens when the angle changes continuously, as it does for a rotating
wheel or an orbiting planet? What path does the corresponding point trace
out over time, and how quickly does it move along that path? The next
chapter begins the study of trigonometry as a theory of motion, taking up
arc length, angular velocity, and the geometry of continuous rotation.

\section{Exercises}
\begin{exercise}
Find the reference angle for each of the following, and state the quadrant
each lies in: $150^\circ$, $-60^\circ$, $300^\circ$, $\dfrac{5\pi}{4}$.
\end{exercise}

\begin{exercise}
Evaluate $\tan(225^\circ)$ and $\sec(150^\circ)$ using reference angles.
\end{exercise}

\begin{exercise}
List three angles coterminal with $40^\circ$, including at least one
negative angle, and explain why all four name the same point of the unit
circle.
\end{exercise}

\begin{exercise}
Explain why $\sin\theta$ and $\cos\theta$ are defined for every real number
$\theta$, while $\tan\theta$ is not.
\end{exercise}
